% Kamp III.2, first supporting lemma: basic past and future patterns.

\paragraph{The first basic-pattern lemma for Theorem III.2.}
Kamp's Theorem~III.2 begins with special finite interval patterns.  Their
existential past and future closures can already be expressed by the
strict operators \(S,U\) over \(\mathsf{Lin}\).  This is Lemma~1 of the
source's proof of III.2; its statement and proof are on printed
pp.~69--70, following the basic-pattern definitions on pp.~63--66
\cite[ch.~III, pp.~60--70]{Kamp1968}.

\begin{hbdefinition}{def:p09-kamp-basic-expressibility}{Expression at a named current point}
Keep the language \(L\), temporal language \(\mathsf{TL}_1\), and map
\(D\) fixed above.  If \(R\) is the unary predicate letter \(P_i\),
write \(\widehat R=q_i\) for its corresponding temporal proposition
letter.  For an \(L\)-formula \(\psi\) with exactly the free variable
\(w\), a temporal formula \(\eta\) \emph{expresses \(\psi\) over
\(\mathsf{Lin}\)} when
\[
 \mathsf{Lin}\models
 \psi\longleftrightarrow D(\eta)[w/t_0].
\]
Rename bound variables before replacing the free current-point variable
\(t_0\) by \(w\).  Equivalence includes every interpretation of the
unary predicates and every assignment to \(w\).  This is the
named-variable form of the source's expressibility convention on p.~59.
It uses Kamp's strict argument order: \(S(A,B)\) or \(U(A,B)\) places
\(A\) at its witness and \(B\) throughout the intervening open interval.
\end{hbdefinition}

\begin{hbdefinition}{def:p09-kamp-basic-classes}{Basic past and future interval patterns}
Let \(x,y\) be distinct variable names.  All bound variables below are
fresh, and all substitutions of endpoint names concern free occurrences.
The letters \(A,B,C\) range over unary predicate letters of \(L\);
they need not be different.

The first binary stages are identical:
\[
\begin{split}
 \operatorname{Bp}_1(x,y)=\operatorname{Bf}_1(x,y)
 =\bigl\{\,&A(x)\land
       \forall v\bigl((x<v\land v<y)\longrightarrow B(v)\bigr)
       \land C(y):\\[-2pt]
       &A,B,C\text{ are unary predicate letters}\,\bigr\}.
\end{split}
\]
For \(n\geq1\), \(\operatorname{Bp}_{n+1}(x,y)\) consists of the
formulas
\[
\begin{gathered}
 A(x)\land\exists z\Bigl(
 x<z\land z<y\land
 \forall v\bigl((x<v\land v<z)\longrightarrow B(v)\bigr)
 \land\phi\Bigr),
 \\
 \phi\in\operatorname{Bp}_n(z,y),
\end{gathered}
\]
and \(\operatorname{Bf}_{n+1}(x,y)\) consists of the formulas
\[
\begin{gathered}
 \exists z\Bigl(
 x<z\land z<y\land\phi\land
 \forall v\bigl((z<v\land v<y)\longrightarrow B(v)\bigr)
 \Bigr)\land A(y),
 \\
 \phi\in\operatorname{Bf}_n(x,z).
\end{gathered}
\]
Each binary formula has exactly its two displayed endpoints free.
The subscripts count the open interval pieces of the pattern.  These are
exact stages, without a clause retaining all formulas of earlier stages.
Put
\[
 \operatorname{Bp}(x,y)=\bigcup_{n\geq1}\operatorname{Bp}_n(x,y),
 \qquad
 \operatorname{Bf}(x,y)=\bigcup_{n\geq1}\operatorname{Bf}_n(x,y).
\]

The corresponding one-variable classes are
\[
\begin{split}
 \operatorname{Bp}_n(y)
   &=\{\exists x\,(x<y\land\phi):
              x\ne y,\ \phi\in\operatorname{Bp}_n(x,y)\},\\
 \operatorname{Bf}_n(x)
   &=\{\exists y\,(x<y\land\phi):
              y\ne x,\ \phi\in\operatorname{Bf}_n(x,y)\},
\end{split}
\]
where \(x\ne y\) here says that the variable names are distinct.
Let \(\operatorname{Bp}(y)\) and \(\operatorname{Bf}(x)\) be their
unions over \(n\geq1\).  The notations
\((\exists x<y)\phi\) and \((\exists y>x)\phi\) abbreviate the two
displayed guarded quantifiers.

The binary base formula does not itself assert \(x<y\).  That guard
is supplied in the one-variable closures.  In particular, an empty open
interval does not remove the requirement for two strictly ordered
endpoint witnesses.  These \(\operatorname{Bp},\operatorname{Bf}\)
classes are distinct from the earlier
\(\operatorname{Bef},\operatorname{Aft},\operatorname{Bet}\) classes
\cite[ch.~III, pp.~63--66]{Kamp1968}.
\end{hbdefinition}

\begin{hbdefinition}{def:p09-kamp-basic-word}{The predicate word and its realizations}
For \(\phi\in\operatorname{Bp}_n(x,y)\cup
\operatorname{Bf}_n(x,y)\), list the unary predicate occurrences from
left to right, including repeated occurrences.  Denote the resulting word by
\[
 \sigma(\phi)=(A_0,B_1,A_1,\ldots,B_n,A_n).
\]
It has \(2n+1\) positions.  Kamp denotes this occurrence sequence by
\(S(\phi)\); \(\sigma\) here distinguishes it from the temporal
operator.  The binary base has word \((A,B,C)\).  A past-stage
constructor prepends \((A,B)\) to its constituent's word; a future-stage
constructor appends \((B,A)\).

In a linear structure \(\mathcal M=\langle T,<,(P_i^{\mathcal M})_i\rangle\),
a \emph{realization} of this word is a sequence of points
\(a_0<\cdots<a_n\) such that
\[
\begin{aligned}
 &a_i\in A_i^{\mathcal M} &&(0\leq i\leq n),\\
 &(a_{i-1},a_i)\subseteq B_i^{\mathcal M} &&(1\leq i\leq n).
\end{aligned}
\]
Here \((a,b)=\{t\in T:a<t<b\}\), and \(R^{\mathcal M}\) is the
interpretation of a unary letter \(R\).  The interval requirement is
vacuous when that interval is empty.  This is the source's sequence
satisfaction relation on pp.~65--66, indexed by \(n\) intervals and
\(n+1\) points.
\end{hbdefinition}

\begin{hblemma}{thm:p09-kamp-iii2-lemma1}{proved=annex:p9.kamp.iii2.lemma1}{Kamp's basic past and future lemma}
Every member of \(\operatorname{Bp}(w)\cup\operatorname{Bf}(w)\)
is expressible over \(\mathsf{Lin}\).

More explicitly, let
\(\sigma=(A_0,B_1,A_1,\ldots,B_n,A_n)\), where \(n\geq1\), be the
word of the binary pattern inside the guarded existential.  Form the
temporal formulas
\[
\begin{aligned}
 L_0&=\widehat A_0,&
 L_i&=\widehat A_i\land S(L_{i-1},\widehat B_i)
            &&(1\leq i\leq n),\\
 R_n&=\widehat A_n,&
 R_i&=\widehat A_i\land U(R_{i+1},\widehat B_{i+1})
            &&(0\leq i<n).
\end{aligned}
\]
Construct the \(L_i\) in increasing index order and the \(R_i\) in
decreasing index order.  If \(\phi\in\operatorname{Bp}_n(x,y)\)
and \(\psi\in\operatorname{Bf}_n(x,y)\) have word \(\sigma\), then
\[
\begin{aligned}
 \mathsf{Lin}\models
  \exists x\,(x<y\land\phi)
       &\longleftrightarrow D(L_n)[y/t_0],\\
 \mathsf{Lin}\models
  \exists y\,(x<y\land\psi)
       &\longleftrightarrow D(R_0)[x/t_0].
\end{aligned}
\]
These are finite constructions using only the input predicate letters and
their corresponding proposition letters.  The seeds \(L_0,R_n\) are
single proposition letters; they do not introduce a stage zero into the
source's basic-pattern classes.  No density, endpointlessness or order
completeness is required \cite[ch.~III, lemma~1, pp.~69--70]{Kamp1968}.
\end{hblemma}

\paragraph{Two interval pieces.}
For a pattern with word \((A_0,B_1,A_1,B_2,A_2)\), its past closure is
expressed by
\[
 \widehat A_2\land
 S\bigl(\widehat A_1\land S(\widehat A_0,\widehat B_1),
         \widehat B_2\bigr),
\]
and its future closure by
\[
 \widehat A_0\land
 U\bigl(\widehat A_1\land U(\widehat A_2,\widehat B_2),
         \widehat B_1\bigr).
\]
Each requires three strictly increasing points.  On a two-point order both
are false, even if every displayed predicate holds everywhere.  For a
one-piece pattern the past formula is false at a first point and the future
formula is false at a last point.  At the other endpoint of a two-point
order, the one-piece formula can be true: its sole open interval is empty,
so its interval predicate imposes no condition.

\paragraph{Role in the remaining construction.}
Theorem~\ref{thm:p09-kamp-iii1} reduces a one-variable formula to Boolean
combinations of current-point and regional components.  The lemma above
supplies the temporal expressions for the basic finite-chain patterns used
in Kamp's subsequent analysis of those components.  The source next treats
the limiting past and future classes in Lemma~2, then develops the remaining
interval lemmas.  Expressibility of all
\(\operatorname{Bef},\operatorname{Aft}\) components over
\(\mathsf{Com}\), and hence Theorem~III.2 itself, is a further result
\cite[ch.~III, pp.~60--62, 66--70]{Kamp1968}.
